\documentclass{article}
\usepackage{/globals/anki}
\ankiprefix{netzwerke}
\ankitopic{Netzwerke}
\ankishow
\begin{document}
\section{Netzwerke}
Mit der richtigen Hardware (Netzkarten, Switches, Router, Leitungen) und der richtigen Software die die relevanten Protokolle implementiert (Treiber, TCP/IP, UDP, HTTP) können mehrere Computer verbunden werden um Daten miteinander auszutauschen.

\subsection{Netzwerkarten}
Es kann zwischen zwei Arten an Netzwerken unterschieden.
\begin{description}
 \item[Local Area Network] Ein LAN ist ein lokaleres Netzwerk, innerhalb eines Hauses, eines Büros oder auch eines ganzen Geländes.

Dies kann mit Kabeln, oder Kabellos (wireless, für ein WLAN) passieren.
 \item[Wide Area Network] Ein WAN ist ein geographisch weiter auseinanderliegendes Netzwerk, wie das gesamte Internet selbst.
\end{description}
\ankinote{arten}{Arten}{LAN, WAN}
\ankinote{lan}{LAN}{Local Area Network; in Haus, Büro, etc}
\ankinote{wan}{WAN}{Wide Area Network; wie Internet}

\subsection{Standards}
Damit die verbundenen Computer sich miteinander auch verstehen können müssen sie die gleiche Sprache sprechen. Dafür werden Standards genutzt auf die sich geeinigt wurde.

Daten werden, vorallem in älteren Systemen, in \emph{XML}, der \emph{Extensible Markup Language}, Ein- und Ausgegeben und übertragen. In der Industrie ist \emph{EDI} (\emph{Electronic Data Interchange}) ein Standard, genutzt werden \emph{Webservices}.
\ankinote{standards}{Standards}{XML, EDI, Webservices}

\end{document}